% =========================================================
% Real-Valued Functions — Notes
% Chapter: functions
% Section: real-valued-functions
% Volume III · Analysis
% =========================================================

\section{Real-Valued Functions}
\label{sec:real-valued-functions}

\begin{tcolorbox}[title={Toolkit --- Real-Valued Functions}]
\begin{itemize}
  \item Pointwise operations: sum, difference, product, quotient,
    scalar multiple, absolute value, max, min. All purely algebraic
    --- no limits involved. $fg$ is pointwise product, not composition.
  \item $\operatorname{BoundedOn}(f,A)$: $\exists B>0,\;\forall x\in A:
    |f(x)|\leq B$.
  \item $\operatorname{MonotoneFunction}(f,A)$: increasing or decreasing.
  \item Maximum attained at $x^*$: $f(x^*) = \sup f(A)$, $x^*\in A$.
    Supremum need not be attained.
\end{itemize}
\end{tcolorbox}

% ---------------------------------------------------------
% Definition — Pointwise Operations
% ---------------------------------------------------------

\begin{tcolorbox}[colback=defbox, colframe=defborder]
\begin{definition}[Pointwise Sum of Functions]
\label{def:pointwise-sum-of-functions}
Let $X\subseteq\mathbb{R}$ and let $f,g:X\to\mathbb{R}$.  The
pointwise sum $f+g:X\to\mathbb{R}$ is the function defined by
\[
  (f+g)(x)=f(x)+g(x)
  \qquad (x\in X).
\]
\end{definition}
\end{tcolorbox}

\begin{remark*}[Standard quantified statement]
\[
  \forall x\in X\;\bigl((f+g)(x)=f(x)+g(x)\bigr).
\]
\end{remark*}

\begin{remark*}[Interpretation]
The sum is computed output by output; at each input, add the two function
values at that same input.
\end{remark*}

\NoLocalDependencies

\begin{tcolorbox}[colback=defbox, colframe=defborder]
\begin{definition}[Pointwise Difference of Functions]
\label{def:pointwise-difference-of-functions}
Let $X\subseteq\mathbb{R}$ and let $f,g:X\to\mathbb{R}$.  The
pointwise difference $f-g:X\to\mathbb{R}$ is the function defined by
\[
  (f-g)(x)=f(x)-g(x)
  \qquad (x\in X).
\]
\end{definition}
\end{tcolorbox}

\begin{remark*}[Standard quantified statement]
\[
  \forall x\in X\;\bigl((f-g)(x)=f(x)-g(x)\bigr).
\]
\end{remark*}

\begin{remark*}[Interpretation]
The difference subtracts the output of $g$ from the output of $f$ at each
input separately.
\end{remark*}

\NoLocalDependencies

\begin{tcolorbox}[colback=defbox, colframe=defborder]
\begin{definition}[Pointwise Product of Functions]
\label{def:pointwise-product-of-functions}
Let $X\subseteq\mathbb{R}$ and let $f,g:X\to\mathbb{R}$.  The
pointwise product $fg:X\to\mathbb{R}$ is the function defined by
\[
  (fg)(x)=f(x)g(x)
  \qquad (x\in X).
\]
\end{definition}
\end{tcolorbox}

\begin{remark*}[Standard quantified statement]
\[
  \forall x\in X\;\bigl((fg)(x)=f(x)g(x)\bigr).
\]
\end{remark*}

\begin{remark*}[Interpretation]
The product is not composition: it multiplies the two real outputs at the
same input.
\end{remark*}

\NoLocalDependencies

\begin{tcolorbox}[colback=defbox, colframe=defborder]
\begin{definition}[Pointwise Scalar Multiple of a Function]
\label{def:pointwise-scalar-multiple-of-a-function}
Let $X\subseteq\mathbb{R}$, let $f:X\to\mathbb{R}$, and let
$c\in\mathbb{R}$.  The pointwise scalar multiple $cf:X\to\mathbb{R}$ is
the function defined by
\[
  (cf)(x)=c f(x)
  \qquad (x\in X).
\]
\end{definition}
\end{tcolorbox}

\begin{remark*}[Standard quantified statement]
\[
  \forall x\in X\;\bigl((cf)(x)=c f(x)\bigr).
\]
\end{remark*}

\begin{remark*}[Interpretation]
Scalar multiplication stretches every output of the function by the same
real scalar.
\end{remark*}

\NoLocalDependencies

\begin{tcolorbox}[colback=defbox, colframe=defborder]
\begin{definition}[Pointwise Absolute Value of a Function]
\label{def:pointwise-absolute-value-of-a-function}
Let $X\subseteq\mathbb{R}$ and let $f:X\to\mathbb{R}$.  The pointwise
absolute value $|f|:X\to\mathbb{R}$ is the function defined by
\[
  |f|(x)=|f(x)|
  \qquad (x\in X).
\]
\end{definition}
\end{tcolorbox}

\begin{remark*}[Standard quantified statement]
\[
  \forall x\in X\;\bigl(|f|(x)=|f(x)|\bigr).
\]
\end{remark*}

\begin{remark*}[Interpretation]
The absolute value operation applies the ordinary real absolute value to
each output.
\end{remark*}

\NoLocalDependencies

\begin{tcolorbox}[colback=defbox, colframe=defborder]
\begin{definition}[Pointwise Maximum of Two Functions]
\label{def:pointwise-maximum-of-two-functions}
Let $X\subseteq\mathbb{R}$ and let $f,g:X\to\mathbb{R}$.  The
pointwise maximum $\max(f,g):X\to\mathbb{R}$ is the function defined by
\[
  \max(f,g)(x)=\max\{f(x),g(x)\}
  \qquad (x\in X).
\]
\end{definition}
\end{tcolorbox}

\begin{remark*}[Standard quantified statement]
\[
  \forall x\in X\;\bigl(\max(f,g)(x)=\max\{f(x),g(x)\}\bigr).
\]
\end{remark*}

\begin{remark*}[Interpretation]
At each input, the new function keeps the larger of the two output values.
\end{remark*}

\NoLocalDependencies

\begin{tcolorbox}[colback=defbox, colframe=defborder]
\begin{definition}[Pointwise Quotient of Functions]
\label{def:pointwise-quotient-of-functions}
Let $X\subseteq\mathbb{R}$ and let $f,g:X\to\mathbb{R}$.  Suppose
$g(x)\neq 0$ for every $x\in X$.  The pointwise quotient $f/g:X\to\mathbb{R}$
is the function defined by
\[
  (f/g)(x)=\frac{f(x)}{g(x)}
  \qquad (x\in X).
\]
\end{definition}
\end{tcolorbox}

\begin{remark*}[Standard quantified statement]
\[
  \bigl(\forall x\in X\;(g(x)\neq 0)\bigr)
  \;\Longrightarrow\;
  \forall x\in X\;\bigl((f/g)(x)=f(x)/g(x)\bigr).
\]
\end{remark*}

\begin{remark*}[Interpretation]
The quotient is defined pointwise only when the denominator function never
vanishes on the domain.
\end{remark*}

\NoLocalDependencies

% ---------------------------------------------------------
% Definition — Bounded Functions
% ---------------------------------------------------------

\begin{tcolorbox}[colback=defbox, colframe=defborder]
\begin{definition}[Bounded Above Function]
\label{def:function-bounded-above}
Let $A\subseteq\mathbb{R}$ and let $f:A\to\mathbb{R}$.  The function
$f$ is bounded above on $A$ if there exists $M\in\mathbb{R}$ such that
\[
  f(x)\leq M
  \qquad (x\in A).
\]
\end{definition}
\end{tcolorbox}

\begin{remark*}[Standard quantified statement]
\[
  \exists M\in\mathbb{R}\;\forall x\in A\;\bigl(f(x)\leq M\bigr).
\]
\end{remark*}

\begin{remark*}[Negated quantified statement]
\[
  \forall M\in\mathbb{R}\;\exists x\in A\;\bigl(f(x)>M\bigr).
\]
\end{remark*}

\begin{remark*}[Failure modes]
No proposed real ceiling works: whenever a candidate upper bound is chosen,
some input has an output larger than that candidate.
\end{remark*}

\begin{remark*}[Failure mode decomposition]
\[
  \forall M\in\mathbb{R}\;\exists x\in A\;\bigl(f(x)>M\bigr).
\]
\end{remark*}

\begin{remark*}[Interpretation]
Bounded above means all output values lie below one common real ceiling.
\end{remark*}

\NoLocalDependencies

\begin{tcolorbox}[colback=defbox, colframe=defborder]
\begin{definition}[Bounded Below Function]
\label{def:function-bounded-below}
Let $A\subseteq\mathbb{R}$ and let $f:A\to\mathbb{R}$.  The function
$f$ is bounded below on $A$ if there exists $m\in\mathbb{R}$ such that
\[
  m\leq f(x)
  \qquad (x\in A).
\]
\end{definition}
\end{tcolorbox}

\begin{remark*}[Standard quantified statement]
\[
  \exists m\in\mathbb{R}\;\forall x\in A\;\bigl(m\leq f(x)\bigr).
\]
\end{remark*}

\begin{remark*}[Negated quantified statement]
\[
  \forall m\in\mathbb{R}\;\exists x\in A\;\bigl(f(x)<m\bigr).
\]
\end{remark*}

\begin{remark*}[Failure modes]
No proposed real floor works: whenever a candidate lower bound is chosen,
some input has an output smaller than that candidate.
\end{remark*}

\begin{remark*}[Failure mode decomposition]
\[
  \forall m\in\mathbb{R}\;\exists x\in A\;\bigl(f(x)<m\bigr).
\]
\end{remark*}

\begin{remark*}[Interpretation]
Bounded below means all output values lie above one common real floor.
\end{remark*}

\NoLocalDependencies

\begin{tcolorbox}[colback=defbox, colframe=defborder]
\begin{definition}[Bounded Function]
\label{def:function-bounded}
Let $A\subseteq\mathbb{R}$ and let $f:A\to\mathbb{R}$.  The function
$f$ is bounded on $A$ if there exists $B>0$ such that
\[
  |f(x)|\leq B
  \qquad (x\in A).
\]
\end{definition}
\end{tcolorbox}

\begin{remark*}[Standard quantified statement]
\[
  \exists B>0\;\forall x\in A\;\bigl(|f(x)|\leq B\bigr).
\]
\end{remark*}

\begin{remark*}[Definition predicate reading]
\[
  \operatorname{BoundedOn}(f,A)
  \;\coloneqq\;
  \exists B>0\;\forall x\in A\;\bigl(|f(x)|\leq B\bigr).
\]
\end{remark*}

\begin{remark*}[Negated quantified statement]
\[
  \forall B>0\;\exists x\in A\;\bigl(|f(x)|>B\bigr).
\]
\end{remark*}

\begin{remark*}[Negation predicate reading]
\[
  \neg\operatorname{BoundedOn}(f,A)
  \;\Longleftrightarrow\;
  \forall B>0\;\exists x\in A\;\bigl(|f(x)|>B\bigr).
\]
\end{remark*}

\begin{remark*}[Failure modes]
Every symmetric interval around $0$ fails to contain all output values of
$f$ on $A$.
\end{remark*}

\begin{remark*}[Failure mode decomposition]
\[
  \forall B>0\;\exists x\in A\;
  \bigl(f(x)>B\;\vee\;f(x)<-B\bigr).
\]
\end{remark*}

\begin{remark*}[Interpretation]
A bounded function has all of its values trapped inside some symmetric
interval $[-B,B]$.  Being bounded above alone is not enough; the function
$f(x)=-x$ on $\mathbb{R}$ is bounded above by $0$ but is not bounded.
\end{remark*}

\begin{remark*}[Dependencies]
\begin{itemize}
  \item \hyperref[def:function-bounded-above]{Bounded Above Function}
  \item \hyperref[def:function-bounded-below]{Bounded Below Function}
\end{itemize}
\end{remark*}

% ---------------------------------------------------------
% Definition — Supremum, Infimum, Maximum, Minimum
% ---------------------------------------------------------

\begin{tcolorbox}[colback=defbox, colframe=defborder]
\begin{definition}[Supremum of a Function on a Set]
\label{def:function-supremum-on-set}
Let $A\subseteq\mathbb{R}$ and let $f:A\to\mathbb{R}$.  The supremum of
$f$ on $A$ is
\[
  \sup_{x\in A} f(x) = \sup f(A),
\]
provided the supremum of $f(A)$ exists.
\end{definition}
\end{tcolorbox}

\begin{remark*}[Standard quantified statement]
\[
  s=\sup_{x\in A}f(x)
  \;\Longleftrightarrow\;
  \bigl(\forall x\in A\;(f(x)\leq s)\bigr)
  \;\wedge\;
  \bigl(\forall u\in\mathbb{R}\;((\forall x\in A\;(f(x)\leq u))\Rightarrow s\leq u)\bigr).
\]
\end{remark*}

\begin{remark*}[Interpretation]
The supremum is the least real ceiling for the output set $f(A)$; it need
not be an output value of $f$.
\end{remark*}

\NoLocalDependencies

\begin{tcolorbox}[colback=defbox, colframe=defborder]
\begin{definition}[Infimum of a Function on a Set]
\label{def:function-infimum-on-set}
Let $A\subseteq\mathbb{R}$ and let $f:A\to\mathbb{R}$.  The infimum of
$f$ on $A$ is
\[
  \inf_{x\in A} f(x) = \inf f(A),
\]
provided the infimum of $f(A)$ exists.
\end{definition}
\end{tcolorbox}

\begin{remark*}[Standard quantified statement]
\[
  t=\inf_{x\in A}f(x)
  \;\Longleftrightarrow\;
  \bigl(\forall x\in A\;(t\leq f(x))\bigr)
  \;\wedge\;
  \bigl(\forall \ell\in\mathbb{R}\;((\forall x\in A\;(\ell\leq f(x)))\Rightarrow \ell\leq t)\bigr).
\]
\end{remark*}

\begin{remark*}[Interpretation]
The infimum is the greatest real floor for the output set $f(A)$; it need
not be an output value of $f$.
\end{remark*}

\NoLocalDependencies

\begin{tcolorbox}[colback=defbox, colframe=defborder]
\begin{definition}[Maximum Point of a Function]
\label{def:function-maximum-point}
Let $A\subseteq\mathbb{R}$ and let $f:A\to\mathbb{R}$.  A point
$x^*\in A$ is a maximum point of $f$ on $A$ if
\[
  f(x)\leq f(x^*)
  \qquad (x\in A).
\]
\end{definition}
\end{tcolorbox}

\begin{remark*}[Standard quantified statement]
\[
  x^*\in A\;\wedge\;\forall x\in A\;\bigl(f(x)\leq f(x^*)\bigr).
\]
\end{remark*}

\begin{remark*}[Negated quantified statement]
\[
  x^*\notin A\;\vee\;\exists x\in A\;\bigl(f(x)>f(x^*)\bigr).
\]
\end{remark*}

\begin{remark*}[Failure modes]
A proposed maximum point fails either by not belonging to the domain or by
being beaten by another output value.
\end{remark*}

\begin{remark*}[Failure mode decomposition]
\[
  x^*\notin A
  \;\vee\;
  \exists x\in A\;\bigl(f(x)>f(x^*)\bigr).
\]
\end{remark*}

\begin{remark*}[Interpretation]
A maximum point is an input where the largest output value is attained.
\end{remark*}

\NoLocalDependencies

\begin{tcolorbox}[colback=defbox, colframe=defborder]
\begin{definition}[Minimum Point of a Function]
\label{def:function-minimum-point}
Let $A\subseteq\mathbb{R}$ and let $f:A\to\mathbb{R}$.  A point
$x_*\in A$ is a minimum point of $f$ on $A$ if
\[
  f(x_*)\leq f(x)
  \qquad (x\in A).
\]
\end{definition}
\end{tcolorbox}

\begin{remark*}[Standard quantified statement]
\[
  x_*\in A\;\wedge\;\forall x\in A\;\bigl(f(x_*)\leq f(x)\bigr).
\]
\end{remark*}

\begin{remark*}[Negated quantified statement]
\[
  x_*\notin A\;\vee\;\exists x\in A\;\bigl(f(x)<f(x_*)\bigr).
\]
\end{remark*}

\begin{remark*}[Failure modes]
A proposed minimum point fails either by not belonging to the domain or by
being undercut by another output value.
\end{remark*}

\begin{remark*}[Failure mode decomposition]
\[
  x_*\notin A
  \;\vee\;
  \exists x\in A\;\bigl(f(x)<f(x_*)\bigr).
\]
\end{remark*}

\begin{remark*}[Interpretation]
A minimum point is an input where the smallest output value is attained.
\end{remark*}

\NoLocalDependencies

% ---------------------------------------------------------
% Definition — Monotone Functions
% ---------------------------------------------------------

\begin{tcolorbox}[colback=defbox, colframe=defborder]
\begin{definition}[Increasing Function]
\label{def:function-increasing}
Let $A\subseteq\mathbb{R}$ and let $f:A\to\mathbb{R}$.  The function
$f$ is increasing on $A$ if
\[
  x_1<x_2\;\Longrightarrow\;f(x_1)\leq f(x_2)
  \qquad (x_1,x_2\in A).
\]
\end{definition}
\end{tcolorbox}

\begin{remark*}[Standard quantified statement]
\[
  \forall x_1,x_2\in A\;\bigl(x_1<x_2\Rightarrow f(x_1)\leq f(x_2)\bigr).
\]
\end{remark*}

\begin{remark*}[Negated quantified statement]
\[
  \exists x_1,x_2\in A\;\bigl(x_1<x_2\;\wedge\;f(x_1)>f(x_2)\bigr).
\]
\end{remark*}

\begin{remark*}[Failure modes]
The function fails to be increasing if some later input has a smaller
output than an earlier input.
\end{remark*}

\begin{remark*}[Failure mode decomposition]
\[
  \exists x_1,x_2\in A\;\bigl(x_1<x_2\;\wedge\;f(x_1)>f(x_2)\bigr).
\]
\end{remark*}

\begin{remark*}[Interpretation]
As the input moves to the right in $A$, the output is not allowed to go
down.
\end{remark*}

\NoLocalDependencies

\begin{tcolorbox}[colback=defbox, colframe=defborder]
\begin{definition}[Strictly Increasing Function]
\label{def:function-strictly-increasing}
Let $A\subseteq\mathbb{R}$ and let $f:A\to\mathbb{R}$.  The function
$f$ is strictly increasing on $A$ if
\[
  x_1<x_2\;\Longrightarrow\;f(x_1)<f(x_2)
  \qquad (x_1,x_2\in A).
\]
\end{definition}
\end{tcolorbox}

\begin{remark*}[Standard quantified statement]
\[
  \forall x_1,x_2\in A\;\bigl(x_1<x_2\Rightarrow f(x_1)<f(x_2)\bigr).
\]
\end{remark*}

\begin{remark*}[Negated quantified statement]
\[
  \exists x_1,x_2\in A\;\bigl(x_1<x_2\;\wedge\;f(x_1)\geq f(x_2)\bigr).
\]
\end{remark*}

\begin{remark*}[Failure modes]
Strict increase fails if some later input has an output that is no larger
than the earlier output.
\end{remark*}

\begin{remark*}[Failure mode decomposition]
\[
  \exists x_1,x_2\in A\;\bigl(x_1<x_2\;\wedge\;f(x_1)\geq f(x_2)\bigr).
\]
\end{remark*}

\begin{remark*}[Interpretation]
Strict increase requires every genuine move to the right in the domain to
produce a genuine increase in output.
\end{remark*}

\begin{remark*}[Dependencies]
\begin{itemize}
  \item \hyperref[def:function-increasing]{Increasing Function}
\end{itemize}
\end{remark*}

\begin{tcolorbox}[colback=defbox, colframe=defborder]
\begin{definition}[Decreasing Function]
\label{def:function-decreasing}
Let $A\subseteq\mathbb{R}$ and let $f:A\to\mathbb{R}$.  The function
$f$ is decreasing on $A$ if
\[
  x_1<x_2\;\Longrightarrow\;f(x_1)\geq f(x_2)
  \qquad (x_1,x_2\in A).
\]
\end{definition}
\end{tcolorbox}

\begin{remark*}[Standard quantified statement]
\[
  \forall x_1,x_2\in A\;\bigl(x_1<x_2\Rightarrow f(x_1)\geq f(x_2)\bigr).
\]
\end{remark*}

\begin{remark*}[Negated quantified statement]
\[
  \exists x_1,x_2\in A\;\bigl(x_1<x_2\;\wedge\;f(x_1)<f(x_2)\bigr).
\]
\end{remark*}

\begin{remark*}[Failure modes]
The function fails to be decreasing if some later input has a larger output
than an earlier input.
\end{remark*}

\begin{remark*}[Failure mode decomposition]
\[
  \exists x_1,x_2\in A\;\bigl(x_1<x_2\;\wedge\;f(x_1)<f(x_2)\bigr).
\]
\end{remark*}

\begin{remark*}[Interpretation]
As the input moves to the right in $A$, the output is not allowed to go up.
\end{remark*}

\NoLocalDependencies

\begin{tcolorbox}[colback=defbox, colframe=defborder]
\begin{definition}[Strictly Decreasing Function]
\label{def:function-strictly-decreasing}
Let $A\subseteq\mathbb{R}$ and let $f:A\to\mathbb{R}$.  The function
$f$ is strictly decreasing on $A$ if
\[
  x_1<x_2\;\Longrightarrow\;f(x_1)>f(x_2)
  \qquad (x_1,x_2\in A).
\]
\end{definition}
\end{tcolorbox}

\begin{remark*}[Standard quantified statement]
\[
  \forall x_1,x_2\in A\;\bigl(x_1<x_2\Rightarrow f(x_1)>f(x_2)\bigr).
\]
\end{remark*}

\begin{remark*}[Negated quantified statement]
\[
  \exists x_1,x_2\in A\;\bigl(x_1<x_2\;\wedge\;f(x_1)\leq f(x_2)\bigr).
\]
\end{remark*}

\begin{remark*}[Failure modes]
Strict decrease fails if some later input has an output that is no smaller
than the earlier output.
\end{remark*}

\begin{remark*}[Failure mode decomposition]
\[
  \exists x_1,x_2\in A\;\bigl(x_1<x_2\;\wedge\;f(x_1)\leq f(x_2)\bigr).
\]
\end{remark*}

\begin{remark*}[Interpretation]
Strict decrease requires every genuine move to the right in the domain to
produce a genuine decrease in output.
\end{remark*}

\begin{remark*}[Dependencies]
\begin{itemize}
  \item \hyperref[def:function-decreasing]{Decreasing Function}
\end{itemize}
\end{remark*}

\begin{tcolorbox}[colback=defbox, colframe=defborder]
\begin{definition}[Monotone Function]
\label{def:function-monotone}
Let $A\subseteq\mathbb{R}$ and let $f:A\to\mathbb{R}$.  The function
$f$ is monotone on $A$ if $f$ is increasing on $A$ or decreasing on $A$.
\end{definition}
\end{tcolorbox}

\begin{remark*}[Standard quantified statement]
\[
  \bigl(\forall x_1,x_2\in A\;(x_1<x_2\Rightarrow f(x_1)\leq f(x_2))\bigr)
  \;\vee\;
  \bigl(\forall x_1,x_2\in A\;(x_1<x_2\Rightarrow f(x_1)\geq f(x_2))\bigr).
\]
\end{remark*}

\begin{remark*}[Definition predicate reading]
\[
  \operatorname{MonotoneFunction}(f,A)
  \;\coloneqq\;
  \bigl(\forall x_1,x_2\in A\;(x_1<x_2\Rightarrow f(x_1)\leq f(x_2))\bigr)
  \;\vee\;
  \bigl(\forall x_1,x_2\in A\;(x_1<x_2\Rightarrow f(x_1)\geq f(x_2))\bigr).
\]
\end{remark*}

\begin{remark*}[Negated quantified statement]
\[
  \bigl(\exists x_1,x_2\in A\;(x_1<x_2\wedge f(x_1)>f(x_2))\bigr)
  \;\wedge\;
  \bigl(\exists y_1,y_2\in A\;(y_1<y_2\wedge f(y_1)<f(y_2))\bigr).
\]
\end{remark*}

\begin{remark*}[Negation predicate reading]
\[
  \neg\operatorname{MonotoneFunction}(f,A)
\]
\end{remark*}

\begin{remark*}[Failure modes]
A nonmonotone function has evidence of both kinds of reversal: one pair
where the outputs go down as inputs increase, and another pair where the
outputs go up as inputs increase.
\end{remark*}

\begin{remark*}[Failure mode decomposition]
\[
  \begin{aligned}
  &\exists x_1,x_2\in A\;(x_1<x_2\wedge f(x_1)>f(x_2))\\
  &\qquad\wedge\;
  \exists y_1,y_2\in A\;(y_1<y_2\wedge f(y_1)<f(y_2)).
  \end{aligned}
\]
\end{remark*}

\begin{remark*}[Interpretation]
A monotone function moves consistently in one direction: nondecreasing or
nonincreasing.
\end{remark*}

\begin{remark*}[Dependencies]
\begin{itemize}
  \item \hyperref[def:function-increasing]{Increasing Function}
  \item \hyperref[def:function-decreasing]{Decreasing Function}
\end{itemize}
\end{remark*}

\begin{tcolorbox}[colback=defbox, colframe=defborder]
\begin{definition}[Constant Function]
\label{def:function-constant}
Let $A\subseteq\mathbb{R}$ and let $f:A\to\mathbb{R}$.  The function
$f$ is constant on $A$ if
\[
  f(x_1)=f(x_2)
  \qquad (x_1,x_2\in A).
\]
\end{definition}
\end{tcolorbox}

\begin{remark*}[Standard quantified statement]
\[
  \forall x_1,x_2\in A\;\bigl(f(x_1)=f(x_2)\bigr).
\]
\end{remark*}

\begin{remark*}[Negated quantified statement]
\[
  \exists x_1,x_2\in A\;\bigl(f(x_1)\neq f(x_2)\bigr).
\]
\end{remark*}

\begin{remark*}[Failure modes]
Constancy fails as soon as two inputs in the domain have different outputs.
\end{remark*}

\begin{remark*}[Failure mode decomposition]
\[
  \exists x_1,x_2\in A\;\bigl(f(x_1)\neq f(x_2)\bigr).
\]
\end{remark*}

\begin{remark*}[Interpretation]
A constant function assigns the same real value to every input in its domain.
\end{remark*}

\NoLocalDependencies
